\documentclass[11pt]{article}
\usepackage{graphicx} % Required for inserting images
\setlength{\parindent}{0pt}
\usepackage[bookmarks=true, colorlinks=true, urlcolor=blue, linkcolor=black]{hyperref}

\usepackage{enumitem}
\usepackage[utf8]{inputenc}
\usepackage[T1]{fontenc}
\usepackage[english]{babel}
\usepackage{geometry}
\usepackage{lastpage}
\usepackage{fancyhdr}
\usepackage{datetime}
\usepackage{parskip}
\usepackage[normalem]{ulem}
\usepackage{xurl}
\urlstyle{same}


\usepackage[absolute]{textpos}
\usepackage{soul}
\usepackage{tikz}
\usetikzlibrary{calc}

% Link to other SOPs, forms
\usepackage{xr}
\externaldocument[102-]{../2_medical/sop_102_medical}
% Definitions for SOP sections 3.
\newcommand{\AdverseEvent}{\textbf{Adverse Event}: Any untoward or unfavorable medical occurrence in a human subject, including any abnormal sign, symptom, or disease, temporally associated with the subject's participation in the research, whether or not considered related to the subject's participation in the research.}

\newcommand{\ControlRoom}{\textbf{Control Room}: B17; The main room of the MRI suite that contains the computer displays and keyboards that control the MRI scanner. The Control Room is accessed from B15 that is controlled by an Electronic Lock.}

\newcommand{\DcmComputer}{\textbf{DICOM Computer}: The Mac Studio located in the Control Room the receives DICOM files from the scanner.}

\newcommand{\ElectronicLock}{\textbf{Electronic Lock}: All doors at the ND-HNC are controlled by electronic locks. The electronic lock system recognizes the owner's Irish1Card, logs the time of the attempted entry, and unlocks the door if the owner has been granted access.}

\newcommand{\EEG}{\textbf{EEG}: Electroencephalograpy, a technique used to measure changes in electromagnetic properties of neural tissue.}

\newcommand{\EquipmentRoom}{\textbf{Equipment Room}: B18; The room containing hardware for MRI operations.}

\newcommand{\FDA}{\textbf{FDA}: Food and Drug Administration.}

\newcommand{\GoodCatch}{\textbf{Good Catch}: Early recognition of an event with the potential to cause damage, injury, or illness.}

\newcommand{\Helper}{\textbf{Helper}: Second person in Zones 3 and 4 who assists in emergency responses, typically Level 2-3.}

\newcommand{\Incident}{\textbf{Incident}: An occurence, action, or event that warrants reporting (e.g. a policy is not followed and/or an individual's health is at risk). Includes Adverse Events, Near Misses, Unanticipated Problems, etc.}

\newcommand{\Lead}{\textbf{Lead}: Primary person in Zones 3 and 4 responsible for the Visitor's safety, in charge during emergency responses. Level 4 operator or MRI Technologist.}

\newcommand{\MagnetRoom}{\textbf{Magnet Room}: B17A; The room within the MRI suite that contains the MRI scanner. It is accessed from within the MRI Control room (B17) through a doorway that is controlled by an Electronic Lock.}

\newcommand{\NearMiss}{\textbf{Near Miss}: An incident that occurred but did not cause personal injury/illness e.g. a small/contained fire, trip on sidewalk or carpeting, falling object.}

\newcommand{\NDHNC}{\textbf{ND-HNC}: Notre Dame Human Neuroimaging Center, Lower Level of Veldman Family Psychology Clinic.}

\newcommand{\Operator}{\textbf{Operator}: MRI Technologist or Level 4 researcher, responsible for conducting research operations and maintaining equpiment. The Lead during emergency responses.}

\newcommand{\PIN}{\textbf{PIN}: Personal Identification Number, separate for Irish1Card and security.}

\newcommand{\ResearchAssistant}{\textbf{Research Assistant}: Individual with access Levels 1-3 to the MRI research suite. Typically the Helper during emergency responses.}

\newcommand{\SeriousInjury}{\textbf{Serious Injury}: An injury or illness that is life threatening; results in permanent impairment of a body function or permanent damage to a body structure; or necessitates medical or surgical intervention to preclude permanent damage or impairment.}

\newcommand{\StimControl}{\textbf{Stimulus Control Computer}: The Windows machines in the Control Room that deploy fMRI tasks and capture participant responses.}

% \newcommand{\Subject}{\textbf{Subject}: An individual who is participating in an experimental protocol at ND-HNC and has not taken the ND-HNC safety course.}

\newcommand{\Supervision}{\textbf{Supervision}: A supervisor is continuously present in both sight and sound range of the individual or task being monitored. This approach ensures the supervisor can actively observe, direct, and immediately intervene if a complication or safety issue arises.}

\newcommand{\TMS}{\textbf{TMS}: Transcranial magnetic stimulation, a technique used to disrupt normal operations of neural tissue.}

\newcommand{\TornadoWarn}{\textbf{Tornado Warning}: A tornado has been sighted or indicated by weather radar and there is imminent danger to life and property. Action is required.}

\newcommand{\TornadoWatch}{\textbf{Tornado Watch}: Tornados are possible in and near the watch area. Be ready to act quickly and move to the designated tornado safe area.}

\newcommand{\UnanticipatedProblem}{\textbf{Unanticipated Problem}: An incident, experience, or outcome that meets \textbf{all} of the following criteria: (1) unexpected (in terms of nature, severity, or frequency) given (a) the research procedures that are described in the protocol-related documents, such as the IRB-approved research protocol and informed consent document; and (b) the characteristics of the subject population being studied; (2) related or possibly related to participation in the research; and (3) suggests that the research places subjects or others at a greater risk of harm than was previously known or recognized.}

\newcommand{\VFPC}{\textbf{VFPC}: Veldman Family Psychology Clinic.}

\newcommand{\Visitor}{\textbf{Visitor}: An individual who has not taken the ND-HNC MRI Safety Course but has a legitimate reason for entering the MRI suites. Research participants and their legal guardians are Visitors. Other Visitors may include: individuals providing/receiving a tour of the VFPC and ND-HNC, faculty and staff who have not received MRI Safety Course, and research assistants.}


\usepackage[super]{cite}
\hypersetup{colorlinks=true, citecolor=black}

% figures and tables
\usepackage{graphicx}
\graphicspath{{../../../../media/sops}{../../../../media/nd_brand}}
\usepackage{rotating}
\usepackage{float}
\usepackage{subfig}
\usepackage{caption}
\DeclareCaptionLabelFormat{adja-page}{\hrulefill\\#1 #2 \emph{(previous page)}}

% Branding
\providecommand\fromlogo{nd_logo.png}
\definecolor{NDblue}{RGB}{12,35,64}
\definecolor{NDgold}{RGB}{174,145,66}
\definecolor{NDgray}{RGB}{193,205,221}

% Center info
\providecommand\nameuniversity{University of Notre Dame}
\providecommand\namecenter{Human Neuroimaging Center}
\providecommand\addrcenter{501 N Hill St. South Bend, IN. 46617}
\providecommand\namecollege{College of Arts and Letters}
\providecommand\addrweb{https://neuroimaging.nd.edu}

% Format header
\urlstyle{sf}
\pagestyle{fancy}

\fancyhead{}
\fancyhead[L]{%
    % LOGO
    \begin{textblock*}{2in}[0.3066,0.39](1.0in,0.5in)
        \includegraphics[width=2in]{\fromlogo}
    \end{textblock*}

    % Above line
    \begin{textblock*}{6.375in}(1.5in,0.32in)
        \sffamily
        \hfill \textcolor{NDgold} \namecenter\ | \textcolor{NDgold}\namecollege
    \end{textblock*}

    % Below line
    \begin{textblock*}{6.375in}(1.5in,0.53in)
        \sffamily
        \hfill \textcolor{NDgold} \addrcenter\ | \textcolor{NDgold}\addrweb
    \end{textblock*}

    %Line
    \begin{tikzpicture}[remember picture,overlay]
        \draw[color=NDblue,line width=0.7pt] (current page.north west)+(2.45in,-0.48in) -- ($(-0.625in,-0.48in)+(current page.north east)$);
    \end{tikzpicture}
}
\renewcommand{\headrulewidth}{0pt}

% Format footer
\fancyfoot[L]{v.0.2.0}
\fancyfoot[C]{\thepage}
\fancyfoot[R]{\textbf{ND-HNC SOP: 104}}

\renewcommand{\footrulewidth}{0.5pt}

% TODO: add language delineating research vs clinical data

% start document
\begin{document}

\renewcommand{\labelenumii}{\arabic{enumi}.\arabic{enumii}}
\renewcommand{\labelenumiii}{\arabic{enumi}.\arabic{enumii}.\arabic{enumiii}}
\renewcommand{\labelenumiv}{\arabic{enumi}.\arabic{enumii}.\arabic{enumiii}.\arabic{enumiv}}

\begin{center}
	\textbf{\Large SOP 104: Incidental Findings}\\
	\hrulefill
\end{center}


\section{Summary}
On rare occasions MRI research personnel may identify neurologic idiosyncracies, termed an `incidental finding', in data collected at the Notre Dame Human Neuroimaging Center (ND-HNC). This document establishes procedures for responding to such incidental findings.

\section{Scope}
These policies apply to all investigators, research assistants, and staff who collect, review, or use MRI data collected at the ND-HNC.

\section{Definitions}

\begin{itemize}
    \item \DICOM
    \item \IRB
    \item \Operator
\end{itemize}


\section{Policies \& Procedures}

% Procedure of escalation:

% 1. If researchers identify a potential incidental finding, communicate with principal investigator (PI)
% 2. PI review, coordinate with ND-HNC staff to review data of interest
% 3. If review confirms original concern, data moved to Incidental Finding location for clinical review
% 4. Radiologist will review, email findings to MRI Director.
% 5. If clinical findings, MRI director will contact PI
% 6. PI schedules an interview with participant, details the incidental finding, recommends participant contact their primary care physician

The ND-HNC operates as a research facility focused on executing experimental MRI protocols for the purpose of scientific investigation. To this end, we provide personnel, infrastructure, and expertise that are appropriate to manage such research data. The ND-HNC is not a clinical facility and we are not organized to deploy clinical MRI protocols for the purpose of medical diagnosis. Given the volume of scans processed at the facility, however, research personnel will occasionally encounter idiosyncratic or anomalous anatomy.

Encountering such idiosyncracies introduces a number of responsibilities which the researcher must navigate to fulfill ethical and legal obligations. Namely, research personnel:

\begin{itemize}
    \item do not normally use diagnostic MRI protocols nor hold the necessary qualifications to make medical diagnoses,
    \item must not cause harm to participants by \textit{commission} (e.g. false alarm leading to undue stress or unnecessary medical expenses) or \textit{omission} (e.g. withholding information which would have been medically relevant),
    \item are required to protect participant anonymity, and
    \item must adhere to IRB regulations of data usage and governance.
\end{itemize}

To help research personnel navigate these obligations, we provide an escalation protocol to be followed whenever incidental findings are detected. A description of these procedures should be included in any IRB application involving use of the ND-HNC MRI scanner. In short, the ND-HNC helps the Principal Investigator (PI) seek a recommendation from a qualified clinician.

\subsection{Initial Review}
All MRI data collected at the ND-HNC should be reviewed by their respective research team. This typically happens during acquisition but may also occur at a later stage. In cases of incidental findings:

\begin{itemize}
    \item Research personnel should \textbf{NOT} share their concerns with the individual who contributed MRI data.
    \item The individual is \textbf{NOT} contacted for follow-up scans nor are they informed that an incidental finding has been detected and the escalation process started.
    \item The PI should be contacted about the incidental finding. It is the responsibility of the PI to coordinate with ND-HNC staff for External Review and with the participant for any potential Follow-Up.
\end{itemize}


\subsection{External Review}
The ND-HNC contracts with qualified clinicians [\textbf{TBD}] to review incidental findings. After a PI contacts ND-HNC staff to initiate an External Review:

\begin{itemize}
    \item Participant structural DICOM files are copied to [\textbf{TBD}], an internal, short-term data storage location.
    \item The clinician reviews the DICOM files either on-site or remotely via the University's secure network.
    \item The clinician sends their review and recommendations by email to nd-hnc@nd.edu and to the director of the ND-HNC, Aron Barbey (abarbey@nd.edu).
\end{itemize}


\subsection{Follow-Up}
ND-HNC staff monitor the email account nd-hnc@nd.edu. Once a review is submitted by the clinician, ND-HNC coordinate with the PI.

If a clinical follow-up is \textbf{NOT} recommended:

\begin{itemize}
    \item The DICOM files are removed from [\textbf{TBD}].
    \item No further action is necessary.
\end{itemize}


Whereas, if a clinical follow-up \textbf{IS} recommended:

\begin{itemize}
    \item The PI schedules a meeting with the participant, sharing the clinician's review and recommendations.
    \item The DICOM files are stored at [\textbf{TBD}] an internal, long-term data storage location, along with the review notes, in case they are requested in the future.
\end{itemize}

\section{Figures}

None.

\end{document}